\begin{claim}{D.58}{Proved}{Negative torus involution sections and cusp fusion}
Let $b$ be a negative involution of the full
once-punctured-torus loop-outer group of
$F=\widehat F(x,y)$, put $c=[x,y]$, and let
$E_b$ be its centre-free extension.
Write $A_2$ for its action on $H_1(F,\Z_2)$.
The $F$-conjugacy classes of order-two sections and
their full fixed-subgroup cohomology are exactly
\[
\begin{array}{c|c|c}
 A_2\bmod2&
 \text{number of section classes}&
 \{\dim_{\F_2}H^1(C_F(S),\F_2)\}_{[S]}\\ \hline
 I&2&\{1,0\}\\
 \ne I&1&\{0\}.
\end{array}                                             \tag{D.58a}
\]
For every such section,
$N_{E_b}(S)=C_F(S)\times S$.
The cusp normalizer is the full profinite dihedral
group
\[
 N_{E_b}(\overline{\langle c\rangle})
       =\overline{\langle c\rangle}
                       \rtimes\langle v\rangle,\qquad
 v^2=1,\quad vcv^{-1}=c^{-1}.                            \tag{D.58b}
\]
Its two internal conjugacy classes of order-two
sections become the same $F$-conjugacy class in $E_b$.
In particular, for every choice of the cusp reflection
$v$, there exists $g\in F$ with
\[
                  gvg^{-1}=cv,\qquad
                  c=[g,v]=g\,(vgv^{-1})^{-1}.            \tag{D.58c}
\]
The dimensions in (D.58a) describe continuous cohomology
of the full closed fixed subgroups.
\end{claim}
\begin{proof}
\emph{The actual extension and the two integral module cases.}
D.27 supplies an actual order-two representative
$\alpha$ and the exact identification
$E_b=F\rtimes_\alpha\langle s\rangle$.
D.43 gives its global multiplier $\lambda=-1$,
so $A_2^2=I$ and $\det A_2=-1$.
Let $u:E_b\to\F_2$ be the quotient character,
$u(s)=1$.

First suppose $A_2\equiv I\pmod2$.
On the free $\Z_2$-module $H^1(F,\Z_2)$, the induced
involution $D$ also satisfies $D\equiv I\pmod2$.
The endomorphism
\[
                         P=(I-D)/2                     \tag{D.58d}
\]
is integral and satisfies $P^2=P$.
Its image and kernel are direct summands of rank one,
because $D$ has the two eigenvalues $1,-1$ over $\Q_2$.
Thus choose continuous characters
$\varphi_+,\varphi_-:F\to\Z_2$ which form a basis and
satisfy
\[
        \varphi_+\alpha=\varphi_+,\qquad
        \varphi_-\alpha=-\varphi_- .                    \tag{D.58e}
\]
This is a decomposition of the specified cohomology
module. All subsequent maps use these actual
characters; no free generating pair is changed.

Their mod-$2$ reductions extend to characters
$a,b:E_b\to\F_2$ by $a(s)=b(s)=0$.
The Hochschild--Serre spectral sequence has the two
rows $\F_2u^i$ and $\F_2^2u^i$.
The degree-one characters extend, so their
transgression vanishes.
All other possible differentials are its products
with powers of $u$ and also vanish.
The associated graded pieces give, for every $n\geq1$,
\[
 H^n(E_b,\F_2)
     =\langle u^n,u^{n-1}a,u^{n-1}b\rangle_{\F_2}.
                                                        \tag{D.58f}
\]
Multiplication by $u$ is an isomorphism in degrees
at least one, since it is an isomorphism on both
associated graded pieces.

\emph{The two Bocksteins and the actual localized ring.}
Reduce (D.58e) modulo four and use the continuous
one-cochains
\[
 \widetilde a(hs^\epsilon)=\varphi_+(h)\bmod4,\qquad
 \widetilde b(hs^\epsilon)=\varphi_-(h)\bmod4.
\]
The first is a homomorphism, while the second has
coboundary
\[
 \delta\widetilde b(hs^\epsilon,ks^\eta)
        =2\epsilon\varphi_-(k)\pmod4.                   \tag{D.58g}
\]
For the coefficient sequence
$0\to\F_2\xrightarrow{2}\Z/4\to\F_2\to0$,
the Bockstein of a degree-one character is its
cup square. This follows directly by taking its
$\{0,1\}$-valued lift; changing the lift changes
the divided coboundary by a coboundary.
Thus (D.58g) gives
\[
                         a^2=0,\qquad b^2=ub.           \tag{D.58h}
\]

Write $ab=\epsilon ua+\delta ub+\gamma u^2$ using
(D.58f), with coefficients in $\F_2$.
Restriction to the chosen section $\langle s\rangle$
gives $\gamma=0$.
Multiplication by $a$ and (D.58h) then give
\[
           0=a^2b
             =\delta\epsilon\,u^2a+\delta^2u^2b .
\]
The two displayed degree-three classes are independent
by (D.58f), so $\delta=0$.
Consequently the degree-zero localization
$R=(H^*(E_b,\F_2)[u^{-1}])^0$ has the exact presentation
\[
 R=\F_2[A,B]\big/
       (A^2,\ B^2-B,\ AB-\epsilon A),\qquad
 A=a/u,\quad B=b/u,\quad \dim_{\F_2}R=3.                 \tag{D.58i}
\]
For either value of $\epsilon$ the map
\[
 \begin{split}
 R&\longrightarrow\F_2[\nu]/(\nu^2)\times\F_2,\\
 1&\longmapsto(1,1),\\
 A&\longmapsto(\nu,0),\\
 B&\longmapsto(\epsilon,1-\epsilon)
 \end{split}                                           \tag{D.58j}
\]
is a ring isomorphism.
It preserves the stated relations and sends the
three-dimensional basis $1,A,B$ to a basis.
In particular $R$ has exactly two primitive
idempotents, and its nilpotent ideal has dimension one.

\emph{The nonidentity action modulo two.}
Now suppose $A_2\bmod2\ne I$.
The dual two-dimensional module
$M=H^1(F,\F_2)$ is the regular $\F_2[C_2]$-module:
$D-I=D+I$ has rank one, with image equal to kernel.
The periodic cyclic resolution therefore gives
\[
       \dim M^{C_2}=1,\qquad H^i(C_2,M)=0\quad(i>0).
\]
Choose the nonzero invariant character on $F$ and
extend it to $\chi:E_b\to\F_2$ with $\chi(s)=0$.
This again makes the only possible transgression zero.
The two-row spectral sequence now gives
\[
 H^1(E_b,\F_2)=\langle u,\chi\rangle,\qquad
 H^n(E_b,\F_2)=\F_2u^n\quad(n\geq2).                    \tag{D.58k}
\]
Both $u\chi$ and $\chi^2$ restrict to zero on
$\langle s\rangle$, whereas $u^2$ restricts nontrivially.
Thus $u\chi=\chi^2=0$.
In this case the degree-zero localization is
\[
             (H^*(E_b,\F_2)[u^{-1}])^0=\F_2.             \tag{D.58l}
\]

\emph{Full Brown restriction and the section labels.}
Every finite subgroup of $E_b$ injects into $E_b/F=C_2$,
because $F$ is torsion-free.
In particular $E_b$ has no subgroup $C_2\times C_2$.
For a section $S$, conjugation in its normalizer
induces the identity on $S$.
Multiplying a normalizer element by the section
element with its quotient label leaves an element
of $C_S=C_F(S)$. This proves the exact direct product
\[
                         N_{E_b}(S)=C_S\times S.        \tag{D.58m}
\]
The closed subgroup $C_S\leq F$ has
$\operatorname{cd}_2(C_S)\leq1$.

The full-profinite Brown restriction theorem applies
with the open subgroup $F$ of cohomological
$2$-dimension one
\cite[Theorem~2 and \S1.1]{A-KochloukovaZalesskii2007}.
For $n>1$ its restriction map to the product of the
cohomologies of (D.58m) is a dense embedding.
The source is finite by (D.58f) or (D.58k);
therefore that map is an isomorphism onto a finite
product, as in the cited proof.
These are the actual restriction maps and preserve
cup products. Since $u$ restricts trivially on
$C_S$ and nontrivially on $S$, localization gives
the ring isomorphism
\[
 R\ \xrightarrow{\ \operatorname{res}\ }\
   \prod_{[S]}
      \left(\F_2\oplus H^1(C_S,\F_2)u^{-1}\right).
                                                        \tag{D.58n}
\]
The second summand in each factor has square zero,
because $H^2(C_S,\F_2)=0$.
Each factor is local, with that summand as its
nilpotent ideal and with residue field $\F_2$.

If $A_2\bmod2\ne I$, (D.58l) and (D.58n) give
one factor and a zero second summand.
If $A_2\bmod2=I$, (D.58j) gives two factors,
one with a one-dimensional second summand and one
with a zero second summand.
This proves all the numbers in (D.58a).
The classes in this comparison are also precisely
the $F$-conjugacy classes:
for a conjugator $e\in E_b$, multiply it by the
inverse of the unique element of $S$ having the
same quotient label. This replaces $e$ by an
element of $F$ without changing $eSe^{-1}$.

In the identity-modulo-two case the two classes
are distinguished by the value of $b$ on their
nonidentity section element.
Indeed the two residue-field evaluations of
(D.58j) give $B=0$ and $B=1$.
Restriction to a section is exactly that
residue-field evaluation in (D.58n).
Also $a$ vanishes on every section element, since
$a^2=0$ whereas the square of the nonzero
character of $C_2$ is nonzero.

\emph{The exact cusp reflections and their fusion.}
Put $L=\overline{\langle c\rangle}\cong\widehat\Z$.
Preservation of the cusp class and A.7 give
\[
       1\longrightarrow L\longrightarrow N_{E_b}(L)
                     \longrightarrow C_2\longrightarrow1.
                                                        \tag{D.58o}
\]
Choose $v$ over its nonidentity label.
The global multiplier $-1$ gives $vcv^{-1}=c^{-1}$.
Write $v^2=c^\kappa$.
Commutation of $v$ with its own square implies
$c^\kappa=c^{-\kappa}$, hence $2\kappa=0$ in
$\widehat\Z$ and $\kappa=0$.
This proves (D.58b), with the multiplication and
specified label of its finite factor retained.

Every section of this dihedral group is
$\langle c^t v\rangle$, $t\in\widehat\Z$.
Conjugation by $c^z$ changes $t$ to $t+2z$,
and conjugation by $v$ changes it to $-t$.
Its two internal classes are therefore represented
by $v$ and $cv$.
These elements have the same image under every
abelian character of $E_b$, since $c=[x,y]\in[F,F]$.
When there is only one full section class they
are already $F$-conjugate.
When there are two, their equal $b$-values and
the preceding residue-field distinction again
put them in the same $F$-class.
An $F$-conjugator sends the unique nonidentity
element to the unique nonidentity element, giving
the first equation of (D.58c).
Multiplication by $v^{-1}$ gives its other
two exact equations.

The executable
\texttt{code/part\_c/torus\_negative\_involution\_d58.py}
checks the mod-four one-cochains, the cyclic
resolution modules, the two explicit localized
algebra isomorphisms, and geometric word examples
of the cusp equation.
The full-group conclusions above follow from
restriction and the actual normalizer maps.
They assert no full triviality or procyclicity
of the fixed subgroups from their first
cohomology dimensions.
\end{proof}
